% Part: proof-theory
% Chapter: sequent-calculus
% Section: invertibility

\documentclass[../../../include/open-logic-section]{subfiles}

\begin{document}

\olfileid[pt]{pt}{seq}{inv}

\olsection{Invertibilidade das Regras}

\begin{defn}
Uma regra~$R$,
\[
\AxiomC{$S_1$}
\AxiomC{$\dots$}
\AxiomC{$S_n$}
\RightLabel{$R$}
\TrinaryInfC{$S$}
\DisplayProof
\]
é \emph{invertível} em um sistema se, sempre que $\Proves S$, então
$\Proves[h_i] S_i$ para $i = 1$, \dots,~$n$. Ela é
\emph{invertível com preservação de !!{height}} (ou \emph{!!{hp}-invertível})
em um sistema se, sempre que $\Proves_n S$, então $\Proves[n] S_i$ para $i =
1$, \dots,~$n$.
\end{defn}

\begin{lem}\ollabel{lem:G3c-invert}
As regras de~\Log{G3c} para $\lnot$, $\land$, $\lor$ e $\lif$ são
invertíveis com preservação de !!{height}, isto é:
\begin{enumerate}
  \item \RightR{\lnot}: Se $\Log{G3c} \Proves[n] \Gamma \Sequent
  \Delta, \lnot !A$, então $\Log{G3c} \Proves[n] !A, \Gamma \Sequent
  \Delta$.
  \item \LeftR{\lnot}: Se $\Log{G3c} \Proves[n] \Gamma, \lnot !A
  \Sequent \Delta$, então $\Log{G3c} \Proves[n] \Gamma \Sequent \Delta,
  !A$.
  \item \RightR{\land}: Se $\Log{G3c} \Proves[n] \Gamma \Sequent
  \Delta, !A \land !B$, então $\Log{G3c} \Proves[n] \Gamma \Sequent
  \Delta, !A$ e $\Log{G3c} \Proves[n] \Gamma \Sequent
  \Delta, !B$.
  \item \LeftR{\land}: Se $\Log{G3c} \Proves[n] \Gamma, !A \land !B
  \Sequent \Delta$, então $\Log{G3c} \Proves[n] !A, !B, \Gamma \Sequent
  \Delta$.
  \item \RightR{\lor}: Se $\Log{G3c} \Proves[n] \Gamma \Sequent
  \Delta, !A \lor !B$, então $\Log{G3c} \Proves[n] \Gamma \Sequent
  \Delta, !A, !B$.
  \item \LeftR{\lor}: Se $\Log{G3c} \Proves[n] !A \lor !B, \Gamma \Sequent
  \Delta$, então $\Log{G3c} \Proves[n] !A, \Gamma \Sequent
  \Delta$ e $\Log{G3c} \Proves[n] !B, \Gamma \Sequent
  \Delta$.
  \item \RightR{\lif}: Se $\Log{G3c} \Proves[n] \Gamma \Sequent
  \Delta, !A \lif !B$, então $\Log{G3c} \Proves[n] !A, \Gamma \Sequent
  \Delta, !B$.
  \item \LeftR{\lif}: Se $\Log{G3c} \Proves[n] !A \lif !B, \Gamma
  \Sequent \Delta$, então $\Log{G3c} \Proves[n] \Gamma \Sequent \Delta,
  !A$ e $\Log{G3c} \Proves[n] !B, \Gamma \Sequent \Delta$.
\end{enumerate}
\end{lem}

\begin{proof}
Precisamos mostrar que, para qualquer regra~$R$ de~\Log{G3c}, se existe
!!a{proof}~$\pi$ da conclusão~$S$ de~$R$, então também existem
!!{proof}s~$\pi_i$ das premissas $S_i$ de~$R$, com $\pheight{\pi_i}
\le \pheight{\pi}$. Considere a regra \LeftR{\land}: suponha que exista
!!a{proof}~$\pi$ com sequente final da forma $!A \land !B, \Gamma
\Sequent \Delta$. Precisamos mostrar que existe !!a{proof} $\pi'$ de
$!A, !B, \Gamma \Sequent \Delta$ com $\pheight{\pi'} \le
\pheight{\pi}$. Fazemos isso por indução sobre $n = \pheight{\pi}$.

Se $n = 0$, $\pi$ consiste apenas em um axioma: $!A \land !B, \Gamma
\Sequent \Delta$. Em \Log{G3c}, os axiomas têm a forma $!C, \Gamma
\Sequent \Delta, !C$, com $!C$ atômica. Em outras palavras, $!C$ não pode ser
$!A \land !B$. Se $!A \land !B, \Gamma \Sequent \Delta$ é um axioma,
então alguma $!C$ atômica deve ser !!a{element} tanto de $\Gamma$
quanto de~$\Delta$. Mas então $!A, !B, \Gamma \Sequent \Delta$ também é um
axioma, e encontramos nossa !!{proof}~$\pi'$ (que também tem
!!{height}~$0$).

Se $n>0$, ela termina em uma inferência. Precisamos provar a afirmação para~$n$,
mas podemos supor que ela é verdadeira para qualquer !!{proof} de altura~$<n$, mesmo
que o contexto seja diferente. Em outras palavras, para quaisquer $k<n$,
$\Pi$ e $\Lambda$, se $!A \land !B, \Pi \Sequent \Lambda$ tem
!!a{proof} de !!{height}~$k$, então existe !!a{proof} de
!!{height}~$\le k$ de $!A, !B, \Pi \Sequent \Lambda$.

Temos dois casos: ou $!A \land
!B$ à esquerda é principal, ou não é. Se for principal, a
última inferência é \LeftR{\land} e a
sub-!!{proof} imediata~$\pi_1$ termina em $!A, !B, \Gamma \Sequent \Delta$. O
resultado vale nesse caso, pois $\pheight{\pi_1} < \pheight{\pi}$.

Se $!A \land !B$ não for principal, alguma outra !!{formula} será, e $!A
\land !B$ pertence ao contexto da última inferência. Então a
hipótese indutiva se aplica às sub-!!{proof}s imediatas da
última regra~$R$, produzindo !!a{proof}~$\pi_1'$ ou duas
!!{proof}s~$\pi_1'$ e~$\pi_2'$ de !!{height} não maior que a
altura das sub-!!{proof}s imediatas de~$\pi$. Os sequentes finais
de $\pi_1'$ e $\pi_2'$ contêm $!A, !B$ no contexto à esquerda
em vez de $!A \land !B$. Aplique a regra~$R$ para obter a
prova~$\pi'$. Por exemplo, suponha que $\pi$ termine em~\LeftR{\lif}:
\[
\AxiomC{}
\RightLabel{$\pi_1$}
\Deduce$!A \land !B, \Gamma' \fCenter \Delta, !C$
\AxiomC{}
\RightLabel{$\pi_2$}
\Deduce$!A \land !B, \Gamma', !D \fCenter \Delta$
\RightLabel{\LeftR{\lif}}
\BinaryInf$!A \land !B, \Gamma', !C \lif !D \fCenter \Delta$
\DisplayProof
\]
Aqui, $!C \lif !D$ à esquerda é principal, e o contexto à
esquerda é $!A \land !B, \Gamma'$ (isto é, $\Gamma = \Gamma', !C \lif !D$).
A hipótese indutiva produz !!{proof}s $\pi_1'$ e $\pi_2'$ de
$!A \land !B, \Gamma' \Sequent \Delta, !C$ e $!A \land !B, \Gamma',
!D \Sequent \Delta$, respectivamente. Esses dois sequentes podem novamente
servir como premissas da regra \LeftR{\lif} para obter~$\pi'$:
\[
\AxiomC{}
\RightLabel{$\pi_1'$}
\Deduce$!A, !B, \Gamma' \fCenter \Delta, !C$
\AxiomC{}
\RightLabel{$\pi_2'$}
\Deduce$!A, !B, \Gamma', !D \fCenter \Delta$
\RightLabel{\LeftR{\lif}}
\BinaryInf$!A, !B, \Gamma', !C \lif !D \fCenter \Delta$
\DisplayProof
\]
Temos
\begin{align*}
\pheight{\pi} & = \max(\pheight{\pi_1}, \pheight{\pi_2}) + 1\\
\pheight{\pi'} & = \max(\pheight{\pi_1'}, \pheight{\pi_2'}) + 1
\intertext{por definição, e}
\pheight{\pi_1'} & \le \pheight{\pi_1}\\
\pheight{\pi_2'} & \le \pheight{\pi_2}
\intertext{pela hipótese indutiva. Portanto,}
\pheight{\pi'} & \le \pheight{\pi}.
\end{align*}
\end{proof}

\begin{prob}
Dê a prova de \olref{lem:G3c-invert} para \RightR{\lif}.
\end{prob}

\begin{lem}\ollabel{lem:invert-quant}
As regras \RightR{\lforall} e \LeftR{\lexists} são
invertíveis com preservação de !!{height} no seguinte sentido:
\begin{enumerate}
  \item Se $\Log{G3c} \Proves[n] \Gamma \Sequent \Delta,
  \lforall[x][!A(x)]$, então $\Log{G3c} \Proves[n] \Gamma \Sequent
  \Delta, !A(c)$; e
  \item se $\Log{G3c} \Proves[n] \lexists[x][!A(x)], \Gamma \Sequent
  \Delta$, então $\Log{G3c} \Proves[n] !A(c), \Gamma \Sequent \Delta$,
\end{enumerate}
para qualquer $c$ que não ocorra em $\Gamma$, $\Delta$ e $\lforall[x][!A(x)]$
ou $\lexists[x][!A(x)]$, respectivamente.
\end{lem}

\begin{proof}
  Mostramos (1) e deixamos (2) como exercício. O argumento é semelhante à
  prova de \olref{lem:G3c-invert}. No passo indutivo, há
  novamente dois casos. O caso em que $\lforall[x][!A(x)]$ é
  principal é rápido: a premissa da última
  inferência $\RightR{\lforall}$ tem a forma exigida $\Gamma
  \Sequent \Delta, !A(a)$, e a sub-!!{proof} imediata~$\pi_1$
  que termina nela tem $\pheight{\pi_1} < \pheight{\pi}$. Além disso, a
  condição de variável própria garante que $a$ não ocorre em
  $\Gamma$, $\Delta$ nem $\lforall[x][!A(x)]$. Como podemos supor que
  $\pi_1$ é regular, $\Subst{\pi_1}{c}{a}$ é !!a{proof} da mesma
  !!{height} de $\Gamma \Sequent \Delta, !A(c)$ por
  \olref[qua]{lem:subst-regular}.

  No segundo caso, $\lforall[x][!A(x)]$ não é principal na
  última inferência; alguma outra fórmula é. Consideremos novamente um caso
  como exemplo. Suponha que a última inferência seja~\RightR{\land}. Alguma
  !!{formula} em $\Delta$ tem a forma~$!B \land !C$ e é principal.
  A !!{proof}~$\pi$ termina em
  \[
\AxiomC{}
\RightLabel{$\pi_1$}
\Deduce$\Gamma \fCenter \Delta', \lforall[x][!A(x)], !B$
\AxiomC{}
\RightLabel{$\pi_2$}
\Deduce$\Gamma \fCenter \Delta', \lforall[x][!A(x)], !C$
\RightLabel{\RightR{\land}}
\BinaryInf$\Gamma \fCenter \Delta', \lforall[x][!A(x)], !B \land !C$
\DisplayProof
\]
em que $\Delta = \Delta', !B \land !C$. Seja $c$ !!a{constant} que
não aparece em $\Delta$; então é claro que também não ocorre em
$\Delta', !B$ nem em $\Delta', !C$. Assim, a hipótese indutiva se aplica
a $\pi_1$ e~$\pi_2$; temos !!{proof}s~$\pi_1'$ e $\pi_2'$ com
$\pheight{\pi_1'} \le \pheight{\pi_1}$ e $\pheight{\pi_2'} \le
\pheight{\pi_2}$. A !!{proof} exigida~$\pi'$ de $\Gamma \Sequent
\Delta, !A(c)$ é:
\[
\AxiomC{}
\RightLabel{$\pi_1'$}
\Deduce$\Gamma \fCenter \Delta', !A(c), !B$
\AxiomC{}
\RightLabel{$\pi_2'$}
\Deduce$\Gamma \fCenter \Delta', !A(c), !C$
\RightLabel{\RightR{\land}}
\BinaryInf$\Gamma \fCenter \Delta', !A(c), !B \land !C$
\DisplayProof
\]
com $\pheight{\pi'} = \max(\pheight{\pi_1'}, \pheight{\pi_2'})+1 \le \pheight{\pi}$.
\end{proof}

\Olref{lem:G3c-invert} desempenha um papel importante na prova do
teorema de eliminação do corte. Mas também pode ser usado para mostrar o
seguinte:

\begin{prop}\ollabel{prop:G3c-cont-adm}
As regras de contração \LeftR{\Contraction} e \RightR{\Contraction}
de~\Log{G1c} são admissíveis em~\Log{G3c} com preservação de !!{height}.
\end{prop}

\begin{proof}
Apresentamos simultaneamente o argumento para \RightR{\Contraction} e
\LeftR{\Contraction}. Suponha que $\pi$ seja !!a{proof} de
$\Gamma \Sequent \Delta, !A, !A$ ou de $!A, !A, \Gamma \Sequent
\Delta$ em~\Log{G3c}. Precisamos mostrar que também existe uma
\Log{G3c}-!!{proof}~$\pi'$ de $\Gamma \Sequent \Delta, !A$ ou $!A,
\Gamma \Sequent \Delta$, respectivamente, com $\pheight{\pi'} \le
\pheight{\pi}$. Fazemos isso por indução sobre $n = \pheight{\pi}$.

Se $n=0$, $\pi$ é apenas um axioma. Mas, se $\Gamma \Sequent \Delta, !A,
!A$ é um axioma de~\Log{G3c}, então $\Gamma \Sequent \Delta, !A$ também é:
ou alguma fórmula atômica~$!C$ é !!a{element} tanto de $\Gamma$ quanto de~$\Delta$,
ou $!A$ é atômica e !!a{element} de~$\Gamma$. O mesmo raciocínio
se aplica se o sequente final for $!A, !A, \Gamma \Sequent
\Delta$.

Se $n>0$, $\pi$ termina em alguma regra~$R$. Se $!A$ não for principal nessa
última inferência, podemos obter !!a{proof}~$\pi'$ como na prova
de \olref{lem:G3c-invert}, aplicando a hipótese indutiva às
sub-!!{proof}s imediatas de~$\pi$ e aplicando a mesma regra~$R$ às
!!{proof}(s) resultantes. A hipótese indutiva afirma que, se existe
!!a{proof} de !!{height}~$k<n$ de $\Pi \Sequent \Lambda, !D,
!D$ ou de $!D, !D, \Pi \Sequent \Lambda$, então existe !!a{proof} de
!!{height}~$\le k$ de $\Pi \Sequent \Lambda, !D$ ou $!D, \Pi \Sequent
\Lambda$, respectivamente. (Observe que a hipótese indutiva se aplica se
os contextos ou a !!{formula} contraída forem diferentes de $\Gamma$,
$\Delta$ ou~$!A$!)

Por exemplo, suponha que $R$ seja $\RightR{\land}$
e que $\pi$ termine em
\[
\AxiomC{}
\RightLabel{$\pi_1$}
\Deduce$\Gamma \fCenter \Delta', !B, !A, !A$
\AxiomC{}
\RightLabel{$\pi_2$}
\Deduce$\Gamma \fCenter \Delta', !C, !A, !A$
\RightLabel{\RightR{\land}}
\BinaryInf$\Gamma \fCenter \Delta', !B \land !C, !A, !A$
\DisplayProof
\]
em que $\Delta = \Delta', !B \land !C$. A hipótese indutiva produz
!!{proof}s $\pi_1'$ e $\pi_2'$ de $\Gamma \fCenter \Delta', !B, !A$
e $\Gamma \fCenter \Delta', !C, !A$, respectivamente, com
$\pheight{\pi_1'} \le  \pheight{\pi_1}$ e $\pheight{\pi_2'} \le
\pheight{\pi_2}$. A !!{proof}~$\pi'$ é:
\[
\AxiomC{}
\RightLabel{$\pi_1'$}
\Deduce$\Gamma \fCenter \Delta', !B, !A$
\AxiomC{}
\RightLabel{$\pi_2'$}
\Deduce$\Gamma \fCenter \Delta', !C, !A$
\RightLabel{\RightR{\land}}
\BinaryInf$\Gamma \fCenter \Delta', !B \land !C, !A$
\DisplayProof
\]

Se $!A$ for principal na última inferência, aplicamos o lema de
inversão (\olref{lem:G3c-invert}) às sub-!!{proof}s imediatas
de~$\pi$, depois a hipótese indutiva e, por fim, novamente a regra~$R$. Por
exemplo, suponha que $\pi$ prove $\Gamma \Sequent \Delta, !A, !A$, que $!A
\ident (!B \land !C)$ e que seja principal na última inferência. Então $\pi$
tem a seguinte forma:
\[
\AxiomC{}
\RightLabel{$\pi_1$}
\Deduce$\Gamma \fCenter \Delta, !B \land !C, !B$
\AxiomC{}
\RightLabel{$\pi_2$}
\Deduce$\Gamma \fCenter \Delta, !B \land !C, !C$
\RightLabel{\RightR{\land}}
\BinaryInf$\Gamma \fCenter \Delta, !B \land !C, !B \land !C$
\DisplayProof
\]
Aplicando \olref{lem:G3c-invert} a $\pi_1$, obtemos !!a{proof}~$\pi_1'$ de
$\Gamma \Sequent \Delta, !B, !B$. (Também obtemos !!a{proof} de
$\Gamma \Sequent \Delta, !C, !B$, mas não precisamos dela.) De modo semelhante,
aplicando \olref{lem:G3c-invert} a~$\pi_2$, obtemos
!!a{proof}~$\pi_2'$ de $\Gamma \Sequent \Delta, !C, !C$. Como a
inversão de \RightR{\land} preserva a !!{height},
$\pheight{\pi_1'} \le  \pheight{\pi_1} < \pheight{\pi}$ e
$\pheight{\pi_2'} \le  \pheight{\pi_2} < \pheight{\pi}$. Portanto, a
hipótese indutiva se aplica a $\pi_1'$ e $\pi_2'$: existem
!!{proof}s~$\pi_1''$ e $\pi_2''$ de $\Gamma \Sequent \Delta, !B$ e
$\Gamma \Sequent \Delta, !C$, respectivamente, com $\pheight{\pi_1''}
\le \pheight{\pi_1'} \le \pheight{\pi_1}$ e $\pheight{\pi_2''} \le
\pheight{\pi_2'} \le \pheight{\pi_2}$. A !!{proof}~$\pi'$ é então:
\[
\AxiomC{}
\RightLabel{$\pi_1''$}
\Deduce$\Gamma \fCenter \Delta, !B$
\AxiomC{}
\RightLabel{$\pi_2''$}
\Deduce$\Gamma \fCenter \Delta, !C$
\RightLabel{\RightR{\land}}
\BinaryInf$\Gamma \fCenter \Delta, !B \land !C$
\DisplayProof
\]
com $\pheight{\pi'} \le \pheight{\pi}$.

Precisamos ter cuidado especial com as regras dos quantificadores nos casos em que
a !!{formula} é principal.

Suponha que $!A \ident \lforall[x][!B(x)]$ seja principal à direita. Então $\pi$ termina em
\[
\AxiomC{}
\RightLabel{$\pi_1$}
\Deduce$\Gamma \fCenter \Delta, \lforall[x][!B(x)], !B(a)$
\RightLabel{\RightR{\lforall}}
\UnaryInf$\Gamma \fCenter \Delta, \lforall[x][!B(x)], \lforall[x][!B(x)]$
\DisplayProof
\]
Por \olref{lem:invert-quant}, existe !!a{proof}~$\pi_1'$ de
!!{height} $\le \pheight{\pi_1}$ de $\Gamma \fCenter \Delta, !B(c),
!B(a)$ para qualquer~$c$ que não ocorra em $\Gamma \fCenter \Delta,
\lforall[x][!B(x)], !B(a)$. Podemos supor que $\pi_1'$ seja regular; assim,
por \olref[qua]{lem:subst-regular}, a !!{proof} $\Subst{\pi_1'}{a}{c}$
é !!a{proof} de $\Gamma \fCenter \Delta, !B(a),
!B(a)$ também de !!{height} $\le \pheight{\pi_1}$. Agora a hipótese
indutiva se aplica e nos dá !!a{proof}~$\pi_1''$ de $\Gamma
\fCenter \Delta, !B(a)$. Aplique a regra \RightR{\lforall} para
obter~$\pi'$:
\[
\AxiomC{}
\RightLabel{$\pi_1''$}
\Deduce$\Gamma \fCenter \Delta, !B(a)$
\RightLabel{\RightR{\lforall}}
\UnaryInf$\Gamma \fCenter \Delta, \lforall[x][!B(x)]$
\DisplayProof
\]
com $\pheight{\pi'} \le \pheight{\pi}$.

Agora suponha que $!A \ident \lexists[x][!B(x)]$
seja principal à direita. Então $\pi$ tem a seguinte forma:
\[
\AxiomC{}
\RightLabel{$\pi_1$}
\Deduce$\Gamma \fCenter \Delta, \lexists[x][!B(x)], \lexists[x][!B(x)], !B(t)$
\RightLabel{\RightR{\lexists}}
\UnaryInf$\Gamma \fCenter \Delta, \lexists[x][!B(x)], \lexists[x][!B(x)]$
\DisplayProof
\]
A hipótese indutiva se aplica a~$\pi_1$ e existem
!!a{proof}s~$\pi_1'$ de $\Gamma \Sequent \Delta, \lexists[x][!B(x)],
!B(t)$. (Observe que nenhum lema de inversão é necessário aqui!) A
!!{proof}~$\pi'$ é então:
\[
\AxiomC{}
\RightLabel{$\pi_1'$}
\Deduce$\Gamma \fCenter \Delta,  \lexists[x][!B(x)], !B(t)$
\RightLabel{\RightR{\lexists}}
\UnaryInf$\Gamma \fCenter \Delta, \lexists[x][!B(x)]$
\DisplayProof
\]
com $\pheight{\pi'} \le \pheight{\pi}$.
\end{proof}

\begin{prob}
Esboce a prova de \olref{prop:G3c-cont-adm} para \LeftR{\Contraction}.
Descreva os casos em que $!A \ident (!B \land !C)$ e $!A \ident (!B
\lif !C)$.
\end{prob}

\begin{prob}
Esboce a prova de \olref{prop:G3c-cont-adm} para \LeftR{\Contraction}
nos casos restantes dos quantificadores, isto é, quando $!A \ident
\lforall[x][!B(x)]$ e $!A \ident \lexists[x][!B(x)]$.
\end{prob}


\end{document}
